\section{考察}\label{sec:discussion}

\subsection{何を分けて評価するか}

Iranian Churnの例では、行単位で分けると同じ入力値を持つ保留行が107件あった。グループ分割は、その数を0にする要件を満たした。同じ入力をまとめても、解約のような異なる結果が同じ組に現れ得るため、重複除去とは挙動が異なる。顧客単位の一般化を問うときに、顧客IDが得られるなら、そのIDに基づく分割を優先して考える。数値入力の一致で作るグループは、抽出した特徴量が変わると変化する。入力の重なり指標も、選んだ列と基準を添えて読む。

自転車貸出の観測では、後の期間でも時刻や稼働日の値は繰り返される。今回確保した境界は日付の境界であり、学習に含まれない暦特徴量の組を作ることとは異なる。時間的な一般化を調べる場合、同じ観測日は分けず、評価期間の開始前に得られた特徴量と正解だけを使えることを確認する必要がある。

\subsection{単純ルール、誤りの内訳、採用判断}

自転車貸出では、曜日・時間帯ルールが全保留期間のmacro F1で上回り、モデルは少数側の「通常」をより多く検出した。誤検知も増えたため、どの誤りを減らしたいかを決めてから数値を比較する必要がある。また、月別のクラス構成も異なり、8月は月初の一部を含まない。月のmacro F1を平均した値は、582行をまとめて混同行列を作った全期間のmacro F1と集計の意味が異なる。比較レポートには両方の分母と対象期間を残す。

評価基盤は、入力データと記録された評価条件から求めた順位、スコア、誤りの内訳を提示する。実務上の採用は利用者が決める。たとえば、交差検証でAのスコアが最も高く、利用者がBを採用する場合は、共通の分割と指標で求めたBとAの差を示せる。保留テストで差を求めるには、両者を同じ保留データに適用した評価が必要になる。

実務への適用では、単純ルールの定義と利用できる特徴量を評価前にそろえ、別期間または別集団で確かめる。自転車のルールは最初の保留結果の確認後に選んだ探索的比較である。したがって、将来の運用を決めるには、同じ比較を事前に固定した追加の後続期間が要る。分類の採用価値を議論するには、見逃しと誤検知の費用、確率の校正、運用時の閾値、実際の施策による結果も必要になる。

\subsection{区間の範囲と再現性}

固定モデルの区間は、観測した行またはグループの再標本化が指標に与える変動を示す。行間やグループ間の依存、データの代表性、特徴量選択やモデル選定までやり直した場合の変動は、評価設計を変えて調べる。日付順の比較では、連続する観測日の依存と観測日の欠落を扱える方法が必要である。連続日をまとめるブロック・ブートストラップ\cite{kuensch}を候補とし、定常性を近似できる期間とブロック長への感度を検討する。十分な観測日を確保し、同じ保留対象に対するモデルとルールの差を集計することが今後の課題である。

保存した入力のハッシュ、分割、予測、設定は、結果の再計算と入力改変の検出に使える。WindowsとPythonでの検証、独立した環境での過去の実行照合、保存したExcelの再計算は、実装の動作と実行再現性に関する証拠である。対象業務の分布への一般化は、時点や組織を変えた追加データから検証する。

\section{結論}

\project{}は、点検したCSVから評価条件を明示し、学習側でモデルを選び、保留テストで単純ルールと誤りの内訳を比較して結果を保存する手順を構築した。通信解約では、特徴量が完全に一致する行をまとめる効果を確認した。自転車貸出では、曜日・時間帯の多数決が全期間のmacro F1で、選定モデルが一部の月と「通常」の検出件数で上回った。日付順に分け、分割、単純ルール、誤りを並べて調べたことで、単一の性能値では見えにくい違いを記録できた。

次は、日付順の分類評価に合わせた性能差の区間と、正解ラベルが遅れて確定する用途の分割を整える。続いて、データ量を変えた学習曲線、学習と選定をやり直した場合の変動、BP画面の初回利用者による結果の読み取りを検証する。
